% Part: incompleteness
% Chapter: representability-in-q
% Section: representing-relations

\documentclass[../../../include/open-logic-section]{subfiles}

\begin{document}

\olfileid{inc}{req}{rel}

\olsection{تمثيل العلاقات}

لنبيّن ما يعنيه أن تكون \emph{علاقة} ما قابلة للتمثيل.

\begin{defn}
\ollabel{defn:representing-relations} تكون العلاقة
$R(x_0,\dots,x_k)$ على الأعداد الطبيعية {\em قابلة للتمثيل في
$\Th{Q}$} إذا وُجدت صيغة $!A_R(x_0,\dots,x_k)$ بحيث كلما صدقت
$R(n_0,\dots,n_k)$ أثبتت $\Th{Q}$
$!A_R(\num{n_0},\dots,\num{n_k})$، وكلما كذبت
$R(n_0,\dots,n_k)$ أثبتت $\Th{Q}$ الصيغة
$\lnot !A_R(\num{n_0},\dots,\num{n_k})$.
\end{defn}

\begin{thm}
\ollabel{thm:representing-rels} تكون العلاقة قابلة للتمثيل في
$\Th{Q}$ إذا وفقط إذا كانت قابلة للحساب.
\end{thm}

\begin{proof}
في الاتجاه المباشر، افترض أن الصيغة $!A_R(x_0,\dots,x_k)$ تمثل
$R(x_0,\dots,x_k)$. وهذه خوارزمية لحساب $R$: عند إدخال
$n_0$، \dots،~$n_k$، ابحث بالتوازي عن برهان
$!A_R(\num{n_0},\dots,\num{n_k})$ وعن برهان
$\lnot !A_R(\num{n_0},\dots,\num{n_k})$. وبحسب فرضنا لا بد أن
يعثر البحث على أحدهما؛ فإن عثر على الأول فأجب «نعم»، وإلا فأجب «لا».

وفي الاتجاه الآخر، افترض أن $R(x_0,\dots,x_k)$ قابلة للحساب. يعني
هذا، بالتعريف، أن الدالة $\Char{R}(x_0,\dots,x_k)$ قابلة للحساب.
وبحسب \olref[int]{thm:representable-iff-comp} تمثل $\Char{R}$ صيغة،
ولتكن $!A_{\Char{R}}(x_0,\dots,x_k,y)$. ولتكن
$!A_R(x_0,\dots,x_k)$ هي الصيغة
$!A_{\Char{R}}(x_0,\dots,x_k,\num{1})$. عندئذ، لأي
$n_0$، \dots،~$n_k$، إذا صدقت $R(n_0,\dots,n_k)$، كانت
$\Char{R}(n_0,\dots,n_k)=1$، فتثبت $\Th{Q}$ الصيغة
$!A_{\Char{R}}(\num{n_0},\dots,\num{n_k},\num{1})$، ومن ثم تثبت
$\Th{Q}$ الصيغة $!A_R(\num{n_0},\dots,\num{n_k})$. ومن جهة أخرى، إذا كذبت
$R(n_0,\dots,n_k)$، كانت $\Char{R}(n_0,\dots,n_k)=0$. وهذا يعني أن
$\Th{Q}$ تثبت
\[
\lforall[y][(!A_{\Char{R}}(\num{n_0}, \dots, \num{n_k}, y) \lif y =
  \num{0})].
\]
ولما كانت $\Th{Q}$ تثبت $\eq/[\num{0}][\num{1}]$، فإن $\Th{Q}$ تثبت
$\lnot !A_{\Char{R}}(\num{n_0},\dots,\num{n_k},\num{1})$، ومن ثم
تثبت $\lnot !A_R(\num{n_0},\dots,\num{n_k})$.
\end{proof}

\begin{prob}
بيّن أنه إذا كانت $R$ قابلة للتمثيل في~$\Th{Q}$، كانت
$\Char{R}$ قابلة للتمثيل أيضًا.
\end{prob}

\end{document}
